\section{讨论}
\label{sec:discussion}

\paragraph{局限} \emph{封闭世界。}所有结果都假设候选原因的清单在调查开始前就已给定。真实分诊中没有人提供这份清单，真实原因也常常不在分析员事先能写出的清单里。缺失时，控制器只能给出兜底的 \textit{other}，或者更糟，给出证据最吻合的那个已列出原因：在 sec-triage 中，把真实原因从清单中移除后，联合证据在 48 个回合中有 24 个给出了另一个需处置的原因（RQ5）。“做决定不需要模型”这一结论成立于这个封闭世界；在调查过程中提出候选原因，很可能正是模型能参与决定的地方，而 \hesp{} 不给它这样做的途径，因为清单之外的结论会被拒绝。\emph{独特观测。}在 sec-triage 和 sigma-triage 中，几乎每个原因都会以概率 1 在某个探针上产生其他原因都不会产生的结果，所以一个选对的探针通常就能定案。这正是由预测表驱动的控制器在那里就足够的原因；在原因共享观测的 comp-triage 中，单一特征规则会把 72 个回合中的 52 个升级。计数表与评测回合来自同一个生成器，而验证器的特征恰好是这些表中的独特结果，因此计数表与生成器表一致、控制器的停止规则与验证器一致，都是由构造决定的。\emph{谁出的题。}探针、每个探针可能返回的结果，以及每个原因会产生哪个结果，在 sec-triage 和 comp-triage 中由我们编写，在 sigma-triage 中由一个模型标注。因此，模型读取器是在一组已经包含大部分领域知识的标签之间做选择；部署时需要分析员为每一类告警写出这些内容，就像今天编写处置手册一样。\emph{来自确定性生成器的表。}统计结果次数是部署时可以做到的，但我们统计的是一个对同一原因每次都返回相同结果的生成器：在 sec-triage 中，每个原因和变体用 1、5 或 20 个开发回合统计出的表，与用 100 个统计出的表相比，只在唯一一个结果随机的探针的 4 个格子上不同（共 90 个格子）。让一个探针就能排除四个原因的那些接近零的数值，反映的是生成器，而不是测得的频率。用已结案工单统计出的表不会这么极端，每个案例需要更多探针，具体多多少我们没有测量；\hesp{} 也没有随攻击手法或日志格式变化而重新统计表的机制。\emph{谁编写了场景。}三个场景都是模拟的，以便控制真值和攻击者；真实遥测更嘈杂、规模更大。它们规模很小（八个原因、十二条规则、三个场景），开发回合与评测回合来自同一个生成器，因此我们的结果能推广到这些任务的新回合，而不能推广到新的原因。sigma-triage 减少了我们对“哪些原因相互竞争”的影响，但其结果是一个模型的确定性标注，并且有一条缺少独特结果的规则被替换。在结构化观测下，环境自己的分类器用固定规则把每个响应变成结果；模型的读取能力只在两种格式的原始日志上检验过，其中漂移格式是我们编写的。\emph{给定且互斥的原因。}原因还是互斥的，环境会报告事件何时变化；并发原因和未报告的变化不在我们的范围内。\emph{我们编写的攻击者。} 我们测试了四种注入文本、两种伪造，以及针对读取器的一种变化后的格式、一种声称措辞和一种日志行伪造。控制两个独立来源组的攻击者能击败佐证，而读取器信任的强度取决于它所信任的规则解析器；我们的解析器按构造对所测攻击免疫，因为它读每个探针的第一条记录。\emph{两个小模型。} 提示比较和读取器攻击只在 Qwen2.5-7B 和 Llama-3.1-8B 上运行。\emph{严格的证据标准。} 我们的验证器会拒绝靠组合观测得出的正确结论；\S\ref{sec:policy} 测量了这一代价，但测量它的 comp-triage 场景很少，且由我们编写。

\paragraph{超出威胁模型的攻击} 有三类攻击超出了 \hesp{} 所达到的目标。控制两个独立来源组的攻击者可以用伪造证据满足佐证；提高所需来源组的数量会抬高门槛，但会让更多诚实的良性案例升级，正如 RQ2 对只有一条特征的原因所显示的那样。污染用于计数预测表的开发案例的攻击者，会使之后的每个结论都发生偏移；预测表应当从结案结果不依赖该智能体的案例中计数。最后，攻击者可以抑制一条特征而不是伪造它，例如避开某个探针能检测到的行为；这会把攻击变成一次升级而不是良性结案，因为 \hesp{} 从不接受靠排除得出的结论。

\paragraph{部署} 我们推荐的 \hesp{} 配置是：信息增益选择、结束守卫、确认条件停止、跨来源佐证和恢复规则；在模型解析原始日志的任何地方，再加上规则解析器优先和读取器信任。结论应当连同所引用的证据和日志一起交给分析员，这样残余的攻击只会造成响应延迟，而不是无人值守的结案。有两个输入决定这套配置的效果。来源组必须按部署中真实的数据流声明，因为佐证的独立性不会超过来源组本身。预测表必须来自与生产环境相似的案例；已结案工单只记录分析员选择运行的探针，据此估计的预测表会面临选择偏差，而且确认条件停止还要求每个原因都至少有一个能把它单独识别出来的探针。

\paragraph{模型的角色} 在结构化观测和给定候选清单下，控制器解释了全部可测的增益；强模型曾经表现出的优势来自等待确认特征，而一条规则就能复现它。模型被需要，是为了读取固定解析器处理不了的日志格式，而这正是它危险的地方。因此，本地分诊智能体中需要加固的组件是读取器，而不是规划器。还有两种角色在我们的场景之外，也是我们接下来会研究的地方：提出清单中缺少的候选原因，以及权衡只有组合起来才成立的证据（如 comp-triage）。两者都需要固定程序解决不了的场景，而单独的控制器是任何此类研究都必须超过的基线。

\paragraph{公开安全数据}\label{sec:realdata} 为了在真实数据上检验这些结果，我们考察了最接近的三个公开数据集能否评估一个调查策略——它需要留下多种合理解释的告警、成本与信息量各异的探针，以及客观的真值（附录~\ref{app:public}）。在我们考察的子集中，没有一个可以：ExCyTIn-Bench~\citep{excytin} 的大多数问题直接点名要定位的告警；在 GUIDE~\citep{guide} 中，组织身份携带了等级 1.50 比特中的 1.03 比特；在我们审计的 OTRF~\citep{otrf} 远程执行录制中，一条基于告警文本的规则就能区分 46 个告警中的 44 个。这些子集缺少的，是客观标注下使用相同工具的攻击与合法管理活动的配对。
